\documentclass[10pt,a4paper]{amsart}
\usepackage{amsmath,amssymb,amsthm,graphicx}
\usepackage{verbatim}
\newenvironment{reference}{\comment}{\endcomment}
%\newenvironment{reference}{}{}
\newenvironment{slogan}{\comment}{\endcomment}
\newenvironment{history}{\comment}{\endcomment}
% For commutative diagrams we use Xy-pic
\usepackage[all]{xy}

% We use 2cell for 2-commutative diagrams.
\xyoption{2cell}
% Theorem environments.
%
\theoremstyle{plain}
\newtheorem{theorem}[subsection]{Theorem}
\newtheorem{proposition}[subsection]{Proposition}
\newtheorem{lemma}[subsection]{Lemma}

\theoremstyle{definition}
\newtheorem{definition}[subsection]{Definition}
\newtheorem{example}[subsection]{Example}
\newtheorem{exercise}[subsection]{Exercise}
\newtheorem{situation}[subsection]{Situation}

\theoremstyle{remark}
\newtheorem{remark}[subsection]{Remark}
\newtheorem{remarks}[subsection]{Remarks}

\numberwithin{equation}{subsection}

% Macros
%
\def\lim{\mathop{\mathrm{lim}}\nolimits}
\def\colim{\mathop{\mathrm{colim}}\nolimits}
\def\Spec{\mathop{\mathrm{Spec}}}
\def\Hom{\mathop{\mathrm{Hom}}\nolimits}
\def\Ext{\mathop{\mathrm{Ext}}\nolimits}
\def\SheafHom{\mathop{\mathcal{H}\!\mathit{om}}\nolimits}
\def\SheafExt{\mathop{\mathcal{E}\!\mathit{xt}}\nolimits}
\def\Sch{\mathit{Sch}}
\def\Mor{\mathop{\mathrm{Mor}}\nolimits}
\def\Ob{\mathop{\mathrm{Ob}}\nolimits}
\def\Sh{\mathop{\mathit{Sh}}\nolimits}
\def\NL{\mathop{N\!L}\nolimits}
\def\CH{\mathop{\mathrm{CH}}\nolimits}
\def\proetale{{pro\text{-}\acute{e}tale}}
\def\etale{{\acute{e}tale}}
\def\QCoh{\mathit{QCoh}}
\def\Ker{\mathop{\mathrm{Ker}}}
\def\Im{\mathop{\mathrm{Im}}}
\def\Coker{\mathop{\mathrm{Coker}}}
\def\Coim{\mathop{\mathrm{Coim}}}

% Boxtimes
%
\DeclareMathSymbol{\boxtimes}{\mathbin}{AMSa}{"02}

%
% Macros for moduli stacks/spaces
%
\def\QCohstack{\mathcal{QC}\!\mathit{oh}}
\def\Cohstack{\mathcal{C}\!\mathit{oh}}
\def\Spacesstack{\mathcal{S}\!\mathit{paces}}
\def\Quotfunctor{\mathrm{Quot}}
\def\Hilbfunctor{\mathrm{Hilb}}
\def\Curvesstack{\mathcal{C}\!\mathit{urves}}
\def\Polarizedstack{\mathcal{P}\!\mathit{olarized}}
\def\Complexesstack{\mathcal{C}\!\mathit{omplexes}}
% \Pic is the operator that assigns to X its picard group, usage \Pic(X)
% \Picardstack_{X/B} denotes the Picard stack of X over B
% \Picardfunctor_{X/B} denotes the Picard functor of X over B
\def\Pic{\mathop{\mathrm{Pic}}\nolimits}
\def\Picardstack{\mathcal{P}\!\mathit{ic}}
\def\Picardfunctor{\mathrm{Pic}}
\def\Deformationcategory{\mathcal{D}\!\mathit{ef}}
\usepackage{cleveref}
\crefname{originallemma}{original lemma}{original lemmas}
\Crefname{originallemma}{Original lemma}{Original lemmas}
\makeatletter
\newlabel{lemma-base-change-envelope}{{0GU7}{source}{}{sourcepointer-1}{}}
\newlabel{lemma-base-change-envelope@cref}{{[originallemma][0][]0GU7}{[1][0][]source}{}{}{}}
\makeatother
\crefname{originalreference}{original reference}{original references}
\Crefname{originalreference}{Original reference}{Original references}
\makeatletter
\newlabel{situation-setup}{{02QL}{source}{}{sourcepointer-2}{}}
\newlabel{situation-setup@cref}{{[originalreference][0][]02QL}{[1][0][]source}{}{}{}}
\makeatother
\title{Computing the restriction of the Chern classes of a perfect object after a proper birational modification}
\author{The Stacks Project Authors}
\date{stacks-project:chow.tex:local-snapshot-20261012; GNU FDL 1.2 or later}
\begin{document}
\maketitle
\noindent\textit{Editorial note.} All delivered mathematical text is the author's own, copied from the pinned Stacks Project chapter file chow.tex; no statement or proof step was rewritten, reordered or removed.\par\medskip
\noindent\textit{Editorial note.} Preamble provenance: the pinned chapter file opens with \textbackslash{}input\{preamble\} and carries no preamble text of its own, so the project's single Stacks-wide preamble staged at sources/stacks-chow/preamble.tex is excerpted verbatim as three spans, lines 8-12 (the verbatim package together with the project's own reference, slogan and history display-environment definitions, including the author's commented-out older form of the reference environment), lines 14-18 (Xy-pic setup) and lines 154-217 (the native theorem declarations and mathematical macros). The lines 8-12 span is required because a counted statement or its surrounding text may use those display environments, which the project typesets as comments. Nothing was added, removed or redefined; the wrapper class and the amsmath/amssymb/amsthm/graphicx base packages are supplied by the assembler.\par\medskip
\noindent\textit{Editorial note.} Same-chapter pointers are located at their \textbackslash{}label lines inside the pinned chapter file; cross-chapter pointers are located in the project's own tag list staged in the snapshot (sources/stacks-chow/tags.txt), at the verbatim lines carrying their canonical tags.\par\medskip
\noindent\textit{Allowed standard prerequisites (editorial).}\begin{itemize}
\item Base change of an envelope. Same chapter chow.tex, label lemma-base-change-envelope, invoked at proof line 9379. Shows that the base change of the trivialising envelope along the given morphism is again an envelope, so that the pulled back object is represented by a complex on the base change.
\item The construction of the Chern classes of a perfect object via an envelope. Same chapter chow.tex, the definition of the restriction of the Chern classes used at proof lines 9388-9390 and claimed/consumed by the envelope construction. Identifies the p-th Chern class of the pulled back complex with the restriction of the Chern class of the object along the morphism, which is the last step of the proof.
\end{itemize}\medskip
\ifdefined\hypertarget\hypertarget{sourcepointer-1}{}\fi
\noindent\textit{Source reference.} Original lemma 0GU7: sources/stacks-chow/pinned-original-source.tex, source lines 3580--3580. Navigation pointer for \textbackslash{}ref\{lemma-base-change-envelope\} used inside the counted span (line 9379); the argument that consumes it is retained in the span.\par\medskip
\ifdefined\hypertarget\hypertarget{sourcepointer-2}{}\fi
\noindent\textit{Source reference.} Original reference 02QL: sources/stacks-chow/pinned-original-source.tex, source lines 1372--1372. Navigation pointer for \textbackslash{}ref\{situation-setup\} used inside the counted span (line 9357); the argument that consumes it is retained in the span.\par\medskip
\begin{lemma}
\label{lemma-chern-classes-computed}
Let $(S, \delta)$ be as in Situation \ref{situation-setup}.
Let $X$ be locally of finite type over $S$. Let $E \in D(\mathcal{O}_X)$
be a perfect object. Assume the Chern classes of $E$ are defined.
For $g : W \to X$ locally of finite type with $W$ integral, there exists
a commutative diagram
$$
\xymatrix{
W' \ar[rd]_{g'} \ar[rr]_b & & W \ar[ld]^g \\
& X
}
$$
with $W'$ integral and $b : W' \to W$ proper birational such that $L(g')^*E$
is represented by a bounded complex $\mathcal{E}^\bullet$ of locally free
$\mathcal{O}_{W'}$-modules of constant rank and we have
$res(c_p(E)) = c_p(\mathcal{E}^\bullet)$ in $A^p(W')$.
\end{lemma}

\begin{proof}
Choose an envelope $f : Y \to X$ such that $Lf^*E$ is isomorphic in
$D(\mathcal{O}_Y)$ to a locally bounded complex $\mathcal{E}^\bullet$
of finite locally free $\mathcal{O}_Y$-modules. The base change
$Y \times_X W \to W$ of $f$ is an envelope by
Lemma \ref{lemma-base-change-envelope}. Choose a point
$\xi \in Y \times_X W$ mapping to the generic point of $W$
with the same residue field. Consider the integral closed subscheme
$W' \subset Y \times_X W$ with generic point $\xi$. The restriction
of the projection $Y \times_X W \to W$ to $W'$ is a proper birational
morphism $b : W' \to W$. Set $g' = g \circ b$. Finally, consider the
pullback $(W' \to Y)^*\mathcal{E}^\bullet$. This is a locally bounded
complex of finite locally free modules on $W'$. Since $W'$ is integral
it follows that it is bounded and that the terms have constant rank.
Finally, by construction $(W' \to Y)^*\mathcal{E}^\bullet$
represents $L(g')^*E$ and by construction its $p$th chern class
gives the restriction of $c_p(E)$ by $W' \to X$. This finishes the proof.
\end{proof}

\end{document}
